\documentclass[12pt]{article}
\providecommand{\palabrasmates}{fracción, numerador, denominador, equivalente…}
\usepackage{album}
\providecommand{\dv}{\mathbin{:}}
\unidad{Fracciones}
\subtitulo{qué es una fracción · equivalentes · comparar y ordenar}

\begin{document}
\portada

% ═══════════════════════════════ LOS CROMOS ═══════════════════════════════
% Los de la cocina: cada idea, con la pizza de la app al lado. Como en el
% álbum de Miut, el andamiaje se retira dentro del guiado: del 1 al 7 casi
% toda la frase está escrita; del 9 en adelante, cada vez menos; el 15, el 16
% y el 17 son solo la pregunta y el espacio.

\debe{\item qué dice el número de abajo y qué dice el de arriba \item cómo se llaman \item un dibujo tuyo, con una pizza}
\begin{cromo}{1}{Qué es una fracción}{clase · la cocina: la pizza cortada}
\begin{minipage}[c]{.72\linewidth}\raggedright
\completa En $\dfrac{3}{4}$, el número de abajo es el \h[2.6cm]{denominador}: dice en cuántos trozos \h[1.6cm]{iguales} se corta la pizza.\par\vspace{2mm}
El de arriba es el \h[2.4cm]{numerador}: dice cuántos trozos \h[2.4cm]{te llevas}.
\end{minipage}%
\begin{minipage}[c]{.28\linewidth}\centering\pizza[escala=.8]{4}{0,1,2}\end{minipage}

\vspace{2mm}
\tuejemplo Una fracción tuya, dibujada con su pizza:
\espacio[1.6cm]{$\dfrac{2}{6}$: la pizza cortada en 6 trozos iguales, y me llevo 2.}
\end{cromo}

\debe{\item dónde va el 0 y dónde el 1 \item cómo se marca una fracción en la regla \item una fracción tuya, marcada}
\begin{cromo}{2}{La fracción en la regla}{app · la regla del tique · Ficha 1}
\completa En la regla, el 1 es una pizza \h[1.6cm]{entera}. Para marcar $\dfrac{3}{8}$, se parte el trozo del 0 al 1 en \h[1cm]{8} partes iguales y se cuentan \h[1cm]{3}.

\vspace{2mm}
\centerline{\regla[10cm]{1}{8}{3/8/tomate}}

\vspace{1mm}
\tuejemplo Marca una fracción tuya en una regla dibujada por ti:
\espacio[1.4cm]{$\dfrac{2}{5}$: del 0 al 1 en 5 partes, y cuento 2.}
\end{cromo}

\debe{\item cuándo una fracción es mayor que 1 \item cómo se escribe como pizzas enteras y un trozo \item un ejemplo tuyo, en la regla}
\begin{cromo}{3}{Más de una pizza}{app · Servir (dos cajas) · Ficha 1 reverso}
\begin{minipage}[c]{.6\linewidth}\raggedright
\completa Si el numerador es \h[2cm]{mayor} que el denominador, la fracción es mayor que 1: no cabe en una sola pizza.\par\vspace{2mm}
$\dfrac{5}{4}$ son \h[.8cm]{1} pizza entera y \h[1cm]{$\frac{1}{4}$} más.
\end{minipage}%
\begin{minipage}[c]{.4\linewidth}\centering\pizza[escala=.65]{4}{0,1,2,3}\ \pizza[escala=.65]{4}{0}\end{minipage}

\vspace{2mm}
\tuejemplo Una fracción tuya mayor que 1, como pizzas enteras y un trozo:
\espacio[1.2cm]{$\dfrac{7}{3}$ = 2 pizzas enteras y $\frac{1}{3}$ más.}
\end{cromo}

\debe{\item qué pasa al cortar cada trozo en 2 \item qué se hace arriba y abajo \item un ejemplo tuyo, con la pizza}
\begin{cromo}{4}{Cortar: amplificar}{app · Servir (el cortador) · Ficha 1}
\begin{minipage}[c]{.6\linewidth}\raggedright
\completa Al cortar cada trozo en 2, salen el \h[1.4cm]{doble} de trozos, la \h[1.4cm]{mitad} de grandes.\par\vspace{2mm}
Se \h[2.4cm]{multiplica} arriba y abajo por el \h[2.4cm]{mismo número}. Esto se llama \textbf{amplificar}.\par\vspace{2mm}
La pizza que te llevas es \h[2.2cm]{la misma}.
\end{minipage}%
\begin{minipage}[c]{.4\linewidth}\centering\cortado{3}{4}{2}\end{minipage}

\vspace{2mm}
\tuejemplo Amplifica una fracción tuya:
\espacio[1.2cm]{$\dfrac{2}{3} = \dfrac{2\cdot 4}{3\cdot 4} = \dfrac{8}{12}$}
\end{cromo}

\debe{\item qué pasa al grapar trozos \item qué se hace arriba y abajo \item un ejemplo tuyo, con la pizza}
\begin{cromo}{5}{Grapar: simplificar}{app · Servir (la grapadora) · Ficha 1}
\begin{minipage}[c]{.6\linewidth}\raggedright
\completa Al grapar los trozos de 2 en 2, hay la \h[1.4cm]{mitad} de trozos, el \h[1.4cm]{doble} de grandes.\par\vspace{2mm}
Se \h[2.2cm]{divide} arriba y abajo entre el \h[2.4cm]{mismo número}. Esto se llama \textbf{simplificar}.\par\vspace{2mm}
Solo se puede grapar si el número \h[3.4cm]{divide a los dos}.
\end{minipage}%
\begin{minipage}[c]{.4\linewidth}\centering\grapado{4}{6}{2}\end{minipage}

\vspace{2mm}
\tuejemplo Simplifica una fracción tuya:
\espacio[1.2cm]{$\dfrac{10}{15} = \dfrac{10\dv 5}{15\dv 5} = \dfrac{2}{3}$}
\end{cromo}

\debe{\item qué son dos fracciones equivalentes \item cómo se ven en la regla \item tres fracciones equivalentes tuyas}
\begin{cromo}{6}{Fracciones equivalentes}{app · Servir · Ficha 1}
\completa Dos fracciones son \textbf{equivalentes} si representan la \h[3.2cm]{misma cantidad} de pizza.\par\vspace{2mm}
En la regla caen en el \h[2.6cm]{mismo punto}: $\dfrac{1}{2}$, $\dfrac{2}{4}$ y $\dfrac{4}{8}$ son el mismo punto con distintos \h[2cm]{nombres}.

\vspace{2mm}
{\centering\regla[8cm]{1}{2}{1/2/tomate}\par\vspace{1mm}\regla[8cm]{1}{4}{2/4/azul}\par\vspace{1mm}\regla[8cm]{1}{8}{4/8/albahaca}\par}

\vspace{1mm}
\tuejemplo Tres fracciones equivalentes a una tuya:
\espacio[1.2cm]{$\dfrac{3}{5} = \dfrac{6}{10} = \dfrac{9}{15} = \dfrac{12}{20}$}
\end{cromo}

\debe{\item qué es la fracción irreducible \item cómo sabes que ya no se puede simplificar \item un ejemplo tuyo, paso a paso}
\begin{cromo}{7}{La fracción irreducible}{app · Servir (de una pieza) · Ficha 1 reverso}
\completa Una fracción es \textbf{irreducible} si ya no se puede \h[2.4cm]{simplificar}: arriba y abajo no tienen ningún divisor común más que el \h[1cm]{1}.\par\vspace{2mm}
$\dfrac{12}{16} = \dfrac{6}{8} =$ \h[1.2cm]{$\frac{3}{4}$}, y $\dfrac{3}{4}$ ya es irreducible porque 3 y 4 \h[4.2cm]{no tienen divisores comunes}.

\vspace{2mm}
\tuejemplo Simplifica una fracción tuya hasta la irreducible:
\espacio[1.2cm]{$\dfrac{18}{24} = \dfrac{9}{12} = \dfrac{3}{4}$}
\end{cromo}

\debe{\item cómo se simplifica de un solo golpe \item qué tiene que ver con Miut \item un ejemplo tuyo}
\begin{cromobrillante}{8}{Simplificar de un golpe: el mcd}{Ficha 1 reverso · Las piezas de Miut}
\completa Para llegar a la irreducible de un golpe, se divide arriba y abajo entre el \h[1.4cm]{mcd} de los dos.\par\vspace{2mm}
$\text{mcd}(18, 24) =$ \h[1cm]{6}, así que $\dfrac{18}{24} = \dfrac{18\dv 6}{24\dv 6} =$ \h[1.2cm]{$\frac{3}{4}$}

\vspace{2mm}
\tuejemplo Una fracción tuya, simplificada de un golpe:
\espacio[1.2cm]{$\dfrac{20}{30}$: mcd(20, 30) = 10, y $\dfrac{20\dv 10}{30\dv 10} = \dfrac{2}{3}$.}
\end{cromobrillante}

\debe{\item cómo se comparan fracciones con el mismo denominador \item un ejemplo tuyo, con el signo}
\begin{cromo}{9}{Comparar: el mismo denominador}{app · Comparar · Ficha 2}
\begin{minipage}[c]{.6\linewidth}\raggedright
\completa Si los trozos son \h[1.6cm]{iguales}, basta contar: es mayor la de \h[3cm]{mayor numerador}.\par\vspace{2mm}
\fA{3}{8}\ \h[.8cm]{$<$}\ \fB{5}{8}
\end{minipage}%
\begin{minipage}[c]{.4\linewidth}\centering\pizza[escala=.65]{8}{0,1,2}\ \pizza[escala=.65]{8}{0,1,2,3,4}\end{minipage}

\vspace{2mm}
\tuejemplo Dos fracciones tuyas con el mismo denominador:
\espacio[1.1cm]{$\dfrac{4}{9} < \dfrac{7}{9}$}
\end{cromo}

\debe{\item cómo se comparan fracciones con el mismo numerador \item por qué el denominador grande da trozos pequeños \item un ejemplo tuyo}
\begin{cromo}{10}{Comparar: el mismo numerador}{app · Comparar · Ficha 2}
\begin{minipage}[c]{.6\linewidth}\raggedright
\completa Si se llevan el mismo número de trozos, gana el trozo más \h[1.8cm]{grande}: el de la pizza cortada en \h[1.6cm]{menos} partes.
\end{minipage}%
\begin{minipage}[c]{.4\linewidth}\centering\pizza[escala=.65]{3}{0}\ \pizza[escala=.65]{4}{0}\end{minipage}

\vspace{2mm}
\palabras ¿Por qué $\frac{1}{4}$ es menor que $\frac{1}{3}$, si el 4 es mayor que el 3?
\espacio[1.3cm]{Porque el 4 dice en cuántos trozos se corta: más comensales, menos ración.}
\end{cromo}

\debe{\item cómo se comparan dos fracciones cualesquiera \item qué es el denominador común \item un ejemplo tuyo, completo}
\begin{cromo}{11}{Reducir a común denominador}{app · Comparar · Ficha 2}
\completa Si los trozos no son iguales, se cortan las dos en la \h[3.2cm]{misma secuencia}: eso es buscar un \textbf{denominador común}. Después, se \h[1.6cm]{cuenta}.

\vspace{2mm}
\tuejemplo Compara dos fracciones tuyas con distinto denominador:
\espacio[1.8cm]{$\dfrac{2}{3}$ y $\dfrac{3}{5}$: en quinceavos, $\dfrac{10}{15}$ y $\dfrac{9}{15}$. Así que $\dfrac{2}{3} > \dfrac{3}{5}$.}
\end{cromo}

\debe{\item cómo se ordenan varias fracciones \item un ejemplo tuyo, en la regla}
\begin{cromo}{12}{Ordenar fracciones}{Ficha 2 reverso · la regla}
\completa Para ordenar varias fracciones:
\espacio[1.6cm]{Las paso todas al mismo denominador (o las marco en la regla) y las ordeno por el numerador: la que está más a la izquierda es la menor.}

\tuejemplo Ordena tres fracciones tuyas de menor a mayor:
\espacio[1.4cm]{$\dfrac{1}{2}$, $\dfrac{3}{8}$, $\dfrac{3}{4}$: en octavos, $\dfrac{4}{8}$, $\dfrac{3}{8}$, $\dfrac{6}{8}$. Orden: $\dfrac{3}{8} < \dfrac{1}{2} < \dfrac{3}{4}$.}
\end{cromo}

\debe{\item cómo se comparan con los productos en cruz \item de dónde sale \item un ejemplo tuyo}
\begin{cromo}{13}{Los productos en cruz}{app · Por teléfono · Ficha 3}
\completa Para comparar $\dfrac{a}{b}$ y $\dfrac{c}{d}$ sin cortar pizza, se multiplica el de arriba de cada una por el \h[2.8cm]{de abajo de la otra}:

\vspace{2mm}
\centerline{\cruzllena{7}{8}{6}{7}}

\vspace{1mm}
Como $49 > 48$, \ \fA{7}{8} \h[.8cm]{$>$} \fB{6}{7}

\vspace{2mm}
\tuejemplo Dos fracciones tuyas, con la cruz:
\espacio[1.2cm]{$\dfrac{3}{8}$ y $\dfrac{2}{5}$: $3\cdot 5 = 15$ y $8\cdot 2 = 16$; $15 < 16$, así que $\dfrac{3}{8} < \dfrac{2}{5}$.}
\end{cromo}

\debe{\item cómo se sabe con la cruz si dos fracciones son equivalentes \item un ejemplo que sí y otro que no}
\begin{cromobrillante}{14}{¿Equivalentes? La cruz lo dice}{app · Por teléfono · Ficha 3 reverso}
\completa $\dfrac{a}{b}$ y $\dfrac{c}{d}$ son equivalentes si $a\cdot d$ \h[.8cm]{$=$} $b\cdot c$.\par\vspace{2mm}
$\dfrac{6}{8}$ y $\dfrac{9}{12}$: \ $6\cdot 12 =$ \h[1cm]{72} \ y \ $8\cdot 9 =$ \h[1cm]{72}: \h[2.4cm]{son equivalentes}.

\vspace{2mm}
\tuejemplo Una pareja tuya que lo sea y otra que no:
\espacio[1.3cm]{$\dfrac{4}{6}$ y $\dfrac{6}{9}$: $36 = 36$, sí.\qquad $\dfrac{3}{4}$ y $\dfrac{5}{7}$: $21 \neq 20$, no.}
\end{cromobrillante}

\debe{\item cómo se juntan dos trozos de distinto tamaño \item un ejemplo tuyo, con la pizza}
\begin{cromo}{15}{Juntar dos trozos}{Ficha 1 reverso · cortar y grapar}
\palabras ¿Cómo se juntan medio pizza y un cuarto? ¿Por qué no se suma arriba y abajo?
\espacio[1.8cm]{Primero los hago del mismo tamaño: $\frac{1}{2} = \frac{2}{4}$. Luego cuento los trozos: $\frac{2}{4} + \frac{1}{4} = \frac{3}{4}$. Abajo no se suma, porque los trozos siguen siendo cuartos.}

\tuejemplo Junta dos trozos tuyos:
\espacio[1.4cm]{$\dfrac{1}{3} + \dfrac{1}{6} = \dfrac{2}{6} + \dfrac{1}{6} = \dfrac{3}{6} = \dfrac{1}{2}$}
\end{cromo}

\debe{\item si entre dos fracciones siempre cabe otra \item cómo se encuentra}
\begin{cromobrillante}{16}{Siempre cabe otra}{Ficha 2 · el reto de la cocina}
\palabras ¿Hay siempre una fracción entre dos fracciones distintas? ¿Cómo la encuentras?
\espacio[1.6cm]{Sí: basta cortar más fino. Entre $\frac{3}{4}$ y $\frac{4}{5}$: en cuarentavos son $\frac{30}{40}$ y $\frac{32}{40}$, y en medio está $\frac{31}{40}$.}

\tuejemplo Una fracción entre $\frac{1}{3}$ y $\frac{1}{2}$:
\espacio[1.2cm]{En doceavos: $\frac{4}{12}$ y $\frac{6}{12}$; en medio, $\frac{5}{12}$.}
\end{cromobrillante}

\debe{\item el error que más has cometido en este tema \item cómo lo evitas}
\begin{cromo}{17}{Mi error típico}{todas las fichas}
\palabras El error que más he cometido en esta unidad, y cómo lo evito:
\espacio[2.6cm]{Este cromo es tuyo: no tiene versión oficial. Un ejemplo de cómo se rellena: «Para comparar $\frac{1}{3}$ y $\frac{1}{4}$ miraba el número de abajo y decía que ganaba el 4. Ahora pienso en la pizza: más trozos, más pequeños».}
\end{cromo}

\end{document}
